\documentclass[12pt,a4paper]{amsart}

\usepackage{amsmath,amssymb,amsthm}
\usepackage{mathtools}
\usepackage[ngerman]{babel}
\usepackage{hyperref}
\usepackage{geometry}
\geometry{margin=2.5cm}

% -------------------------------------------------------
% Theorem-Umgebungen
% -------------------------------------------------------
\newtheorem{theorem}{Satz}[section]
\newtheorem{lemma}[theorem]{Lemma}
\newtheorem{corollary}[theorem]{Korollar}
\newtheorem{proposition}[theorem]{Proposition}
\theoremstyle{definition}
\newtheorem{definition}[theorem]{Definition}
\newtheorem{remark}[theorem]{Bemerkung}
\newtheorem{example}[theorem]{Beispiel}

% -------------------------------------------------------
% Eigene Befehle
% -------------------------------------------------------
\newcommand{\N}{\mathbb{N}}
\newcommand{\Z}{\mathbb{Z}}
\newcommand{\Q}{\mathbb{Q}}
\newcommand{\R}{\mathbb{R}}
\newcommand{\C}{\mathbb{C}}
\newcommand{\Fp}{\mathbb{F}_p}
\DeclareMathOperator{\Re}{Re}

% -------------------------------------------------------
\begin{document}
% -------------------------------------------------------

\title{Die \texorpdfstring{$L$}{L}-Funktion einer elliptischen Kurve:\\
       Euler-Produkt, Frobenius-Spur, analytische Fortsetzung\\
       und Funktionalgleichung}

\author{Michael Fuhrmann}
\address{Specialist Mathematics Project, Build 84}
\email{specialist-maths@localhost}
\date{\today}

\begin{abstract}
Die $L$-Funktion $L(E,s)$ einer elliptischen Kurve $E/\Q$ kodiert globale
arithmetische Information durch ihre Lokalfaktoren.
F\"ur jede Primzahl $p$ guter Reduktion wird der Lokalfaktor durch die
\emph{Frobenius-Spur} $a_p = p + 1 - \#E(\Fp)$ bestimmt;
Hasse bewies $|a_p| \le 2\sqrt{p}$.
Die $L$-Funktion ist zun\"achst f\"ur $\Re(s) > 3/2$ durch ein Euler-Produkt
definiert.
Der tiefe \emph{Modularit\"atssatz} (Wiles 1995, Breuil--Conrad--Diamond--Taylor 2001)
beweist die analytische Fortsetzung auf $\C$ und die Funktionalgleichung,
die $L(E,s)$ und $L(E,2-s)$ verbindet.
\end{abstract}

\maketitle
\tableofcontents

% ================================================================
\section{Einleitung}
% ================================================================

Die $L$-Funktion einer elliptischen Kurve wurde von Hasse (1930er) und
Weil (1950er) untersucht, motiviert durch die Analogie zur Riemann'schen
Zeta-Funktion.
Die zentrale lokale Zutat ist die Anzahl der Punkte $\#E(\Fp)$ f\"ur jede
Primzahl $p$.
Die globale $L$-Funktion verpackt diese lokalen Daten in eine einzige
analytische Funktion; die Birch--Swinnerton-Dyer-Vermutung (Paper~27)
sagt vorher, dass ihr Verhalten bei $s=1$ den Rang von $E(\Q)$ bestimmt.

% ================================================================
\section{Punkte modulo Primzahlen z\"ahlen}
% ================================================================

\begin{definition}[Frobenius-Spur]
\label{def:ap}
F\"ur eine elliptische Kurve $E/\Q$ und eine Primzahl $p$ guter Reduktion
ist die \emph{Frobenius-Spur} definiert als
\[
  a_p \;=\; p + 1 - \#E(\Fp).
\]
Dabei z\"ahlt $\#E(\Fp)$ alle affinen Punkte $(x,y) \in \Fp^2$ mit
$y^2 = x^3 + ax + b$ (modulo $p$ reduziert) plus den Punkt im Unendlichen.
\end{definition}

\begin{theorem}[Hasses Satz, 1933]
\label{thm:hasse}
F\"ur jede Primzahl $p$ guter Reduktion gilt
\[
  |a_p| \;\le\; 2\sqrt{p}.
\]
Gleichbedeutend: $\#E(\Fp) = p+1-a_p \in [p+1-2\sqrt{p},\; p+1+2\sqrt{p}]$.
\end{theorem}

\begin{proof}[Beweisskizze]
Schreibe $a_p = \alpha + \bar\alpha$, wo $\alpha, \bar\alpha \in \C$ die
Eigenwerte des Frobenius auf dem $\ell$-adischen Tate-Modul $T_\ell(E)$
sind, mit $\alpha\bar\alpha = p$.
Dann $|a_p| \le 2|\alpha| = 2\sqrt{p}$.
Die Identit\"at $|\alpha|^2 = p$ entspricht der Riemann-Hypothese f\"ur die
Zeta-Funktion von $E/\Fp$, die Hasse direkt f\"ur elliptische Kurven bewies.
Moderner Beweis via Weil-Kohomologie; vgl.\ \cite{Silverman1986}, Anhang~C.
\end{proof}

\begin{example}[Berechnung von $a_p$]
\label{ex:ap}
F\"ur $E: y^2 = x^3 - x$ und $p = 5$:
Affine Punkte modulo~$5$: $(0,0), (1,0), (2,1), (2,4), (3,2), (3,3), (4,0)$
-- also $7$ affine Punkte plus $\mathcal{O}$, ergibt $\#E(\mathbb{F}_5) = 8$.
Damit $a_5 = 5+1-8 = -2$. Und $|-2| \le 2\sqrt{5} \approx 4{,}47$. \checkmark
\end{example}

% ================================================================
\section{Das Euler-Produkt}
% ================================================================

\begin{definition}[Lokale $L$-Faktoren]
\label{def:local_factors}
F\"ur eine Primzahl $p$ ist der \emph{lokale $L$-Faktor} von $E$:
\begin{itemize}
  \item \emph{Gute Reduktion} ($p \nmid N_E$):
    \[
      L_p(E,s)^{-1} = 1 - a_p\,p^{-s} + p^{1-2s}.
    \]
  \item \emph{Multiplikative Reduktion} ($p \| N_E$):
    \[
      L_p(E,s)^{-1} = 1 - \epsilon_p\,p^{-s}, \quad \epsilon_p \in \{+1,-1\}.
    \]
  \item \emph{Additive Reduktion} ($p^2 \mid N_E$):
    \[
      L_p(E,s)^{-1} = 1.
    \]
\end{itemize}
\end{definition}

\begin{definition}[$L$-Funktion von $E$]
\label{def:L_function}
Die \emph{$L$-Funktion} von $E$ ist das Euler-Produkt
\[
  L(E,s) \;=\; \prod_p L_p(E,s)^{-1}
  \;=\; \prod_{p \nmid N_E} \frac{1}{1 - a_p p^{-s} + p^{1-2s}}
        \cdot \prod_{p \mid N_E} \frac{1}{1 - \epsilon_p p^{-s}}.
\]
\end{definition}

\begin{proposition}[Dirichlet-Reihen-Darstellung]
\label{prop:dirichlet}
F\"ur $\Re(s) > 3/2$ konvergiert das Euler-Produkt absolut und ist gleich
\[
  L(E,s) = \sum_{n=1}^{\infty} \frac{a_n}{n^s},
\]
wobei $a_1 = 1$, $a_p = $ Frobenius-Spur f\"ur $p \nmid N_E$,
und h\"ohere Primzahlpotenzen durch die Rekursion
$a_{p^k} = a_p\,a_{p^{k-1}} - p\,a_{p^{k-2}}$ bestimmt werden.
\end{proposition}

% ================================================================
\section{Die vervollst\"andigte $L$-Funktion und Funktionalgleichung}
% ================================================================

\begin{definition}[Vervollst\"andigte $L$-Funktion]
\label{def:completed_L}
Die \emph{vervollst\"andigte $L$-Funktion} ist
\[
  \Lambda(E,s) \;=\; \left(\frac{\sqrt{N_E}}{2\pi}\right)^s
  \Gamma(s)\,L(E,s).
\]
\end{definition}

\begin{theorem}[Funktionalgleichung via Modularit\"at]
\label{thm:functional_eq}
Die vervollst\"andigte $L$-Funktion erf\"ullt
\begin{equation}
  \label{eq:functional}
  \Lambda(E,s) \;=\; \epsilon(E)\,\Lambda(E,2-s),
\end{equation}
wobei $\epsilon(E) \in \{+1,-1\}$ die \emph{Wurzelzahl} von $E$ ist.
\end{theorem}

\begin{proof}[Beweis via Modularit\"at]
Nach dem Modularit\"atssatz
(Wiles \cite{Wiles1995} f\"ur semistabile Kurven,
Breuil--Conrad--Diamond--Taylor \cite{BCDT2001} allgemein)
entspricht $E$ einer Neuform $f = \sum a_n q^n$ vom Gewicht~$2$ und
Niveau~$N_E$.

Die $L$-Funktion einer Neuform vom Gewicht~$2$ erf\"ullt die
Funktionalgleichung nach der klassischen Hecke-Theorie (1936):
$\Lambda(f,s) = \epsilon_f\,\Lambda(f,2-s)$.
Da $L(E,s) = L(f,s)$, folgt die Behauptung.
\end{proof}

\begin{remark}[Wurzelzahl und Parit\"atsvermutung]
Gilt $\epsilon(E) = -1$, so folgt aus der Funktionalgleichung $L(E,1) = 0$.
Die \emph{Parit\"atsvermutung} besagt: $(-1)^r = \epsilon(E)$.
\end{remark}

% ================================================================
\section{Der Modularit\"atssatz}
% ================================================================

\begin{theorem}[Modularit\"atssatz]
\label{thm:modularity}
Jede elliptische Kurve $E/\Q$ ist \emph{modular}: Es gibt eine Neuform
$f \in S_2(\Gamma_0(N_E))$ mit
\[
  L(E,s) = L(f,s) = \sum_{n=1}^{\infty} \frac{a_n(f)}{n^s},
\]
wobei $a_n(f)$ die Fourier-Koeffizienten von $f$ sind.
\"Aquivalent: Es gibt einen nicht-konstanten Morphismus
$\varphi: X_0(N_E) \to E$ algebraischer Kurven \"uber $\Q$.
\end{theorem}

\begin{proof}[Historische Anmerkungen]
Vermutet von Taniyama (1955) und Shimura (1957)
(Shimura--Taniyama-Vermutung, sp\"ater Shimura--Taniyama--Weil).
\begin{itemize}
  \item Wiles \cite{Wiles1995} bewies dies f\"ur semistabile Kurven.
    Der Beweis verwendet Galois-Darstellungen und das Schl\"usselergebnis:
    Eine Surjektion $R \to T$ (Deformationsring auf Hecke-Ring) ist ein
    Isomorphismus.
  \item Breuil, Conrad, Diamond, Taylor \cite{BCDT2001} vervollst\"andigten
    den Beweis f\"ur alle elliptischen Kurven \"uber $\Q$.
\end{itemize}
Fermats Letzter Satz folgt als Korollar (\"uber Ribets Satz: eine
hypothetische L\"osung ergäbe eine Frey-Kurve, die nicht modular w\"are).
\end{proof}

% ================================================================
\section{Die $L$-Funktion bei \texorpdfstring{$s=1$}{s=1}}
% ================================================================

\begin{definition}[Verschwindungsordnung und f\"uhrender Koeffizient]
\label{def:order_vanishing}
Die \emph{analytische Verschwindungsordnung} von $L(E,s)$ bei $s=1$ ist
\[
  r_{\mathrm{an}} \;=\; \ord_{s=1} L(E,s).
\]
Der \emph{f\"uhrende Koeffizient} ist
$L^*(E,1) = \lim_{s \to 1} (s-1)^{-r_{\mathrm{an}}} L(E,s)$.
Die BSD-Vermutung sagt $r_{\mathrm{an}} = r$ (der algebraische Rang) vorher.
\end{definition}

\begin{remark}[Parit\"at aus der Funktionalgleichung]
F\"ur $\epsilon(E) = -1$ gibt die Funktionalgleichung~\eqref{eq:functional}
$L(E,1) = -L(E,1)$, also $L(E,1) = 0$ und $r_{\mathrm{an}} \ge 1$.
\end{remark}

% ================================================================
\section{Schlussfolgerung}
% ================================================================

Die $L$-Funktion einer elliptischen Kurve ist die Br\"ucke zwischen lokalen
arithmetischen Daten (den Frobenius-Spuren $a_p$) und globaler Arithmetik
(dem Rang von $E(\Q)$).
Hasses Satz sichert $|a_p| \le 2\sqrt{p}$ und damit Konvergenz f\"ur $\Re(s) > 3/2$.
Der Modularit\"atssatz (Wiles et al.) liefert die analytische Fortsetzung
und die Funktionalgleichung~\eqref{eq:functional} mit Wurzelzahl $\epsilon(E) = \pm 1$.

Die zentrale Frage -- welche Ordnung hat das Verschwinden $r_{\mathrm{an}}$
bei $s = 1$? -- ist Gegenstand der Birch--Swinnerton-Dyer-Vermutung (Paper~27).

% ================================================================
\begin{thebibliography}{10}

\bibitem{Silverman1986}
J.~H.~Silverman,
\emph{The Arithmetic of Elliptic Curves},
Graduate Texts in Mathematics~106,
Springer, New York, 1986.

\bibitem{Wiles1995}
A.~Wiles,
Modular elliptic curves and Fermat's last theorem,
\emph{Ann.\ of Math.}\ (2) \textbf{142} (1995), 443--551.

\bibitem{BCDT2001}
C.~Breuil, B.~Conrad, F.~Diamond und R.~Taylor,
On the modularity of elliptic curves over $\Q$,
\emph{J.\ Amer.\ Math.\ Soc.}\ \textbf{14} (2001), 843--939.

\bibitem{Cornell1997}
G.~Cornell, J.~H.~Silverman und G.~Stevens (Hrsg.),
\emph{Modular Forms and Fermat's Last Theorem},
Springer, New York, 1997.

\bibitem{Knapp1992}
A.~W.~Knapp,
\emph{Elliptic Curves},
Princeton University Press, 1992.

\bibitem{TaylorWiles1995}
R.~Taylor und A.~Wiles,
Ring-theoretic properties of certain Hecke algebras,
\emph{Ann.\ of Math.}\ (2) \textbf{141} (1995), 553--572.

\end{thebibliography}

\end{document}
